\begin{lemma}[MellinOfDeltaSpikeAt1-asymp]
\label{MellinOfDeltaSpikeAt1_asymp}
\lean{MellinOfDeltaSpikeAt1_asymp}
% at /home/runner/work/trough/trough/library/PrimeNumberTheoremAnd/MellinCalculus.lean:700.0-764.41
\leanok
As $\epsilon\to 0$, we have
  $$\mathcal{M}(\nu_\epsilon)(1) = 1+O(\epsilon).$$
\end{lemma}
\begin{proof}
\leanok
By Lemma \ref{MellinOfDeltaSpikeAt1},
  $$
    \mathcal M(\nu_\epsilon)(1)=\mathcal M(\nu)(\epsilon)
  $$
  which by Definition \ref{MellinTransform} is
  $$
    \mathcal M(\nu)(\epsilon)=\int_0^\infty\nu(x)x^{\epsilon-1}dx
    .
  $$
  Since $\nu(x) x^{\epsilon-1}$ is integrable (because $\nu$ is continuous and compactly
  supported),
  $$
    \mathcal M(\nu)(\epsilon)-\int_0^\infty\nu(x)\frac{dx}x
    =\int_0^\infty\nu(x)(x^{\epsilon-1}-x^{-1})dx
    .
  $$
  By Taylor's theorem,
  $$
    x^{\epsilon-1}-x^{-1}=O(\epsilon)
  $$
  so, since $\nu$ is absolutely integrable,
  $$
    \mathcal M(\nu)(\epsilon)-\int_0^\infty\nu(x)\frac{dx}x=O(\epsilon)
    .
  $$
  We conclude the proof using Theorem \ref{SmoothExistence}.
\end{proof}
